% !TEX spellcheck = en_US



% ~~~~~~~~~~~~~~~~~~~~~~~~~~~~~~~~~~~~~~~ V
% (NC) 2026 Haluk Bingol
% github.com/halukbingol/LaTeX-Templates
%
% Licensees may copy, distribute, display, and perform the work and make derivative works 
% and remixes based on it only for non-commercial purposes. 
% ~~~~~~~~~~~~~~~~~~~~~~~~~~~~~~~~~~~~~~~ A




% =====================================================================
%  Java on the Command Line -- Sophomore Level: Projects, Classpath, JARs and Build Scripts
%  ---------------------------------------------------------------------
%  Built on hbCisLaTeXTemplate.tex with the libraries in ../../hblib/.
%  Programs and their outputs are read from codeoutput/:
%     codeoutput/<Name>.java   the program
%     codeoutput/<Name>.in     keyboard input (optional)
%     codeoutput/<Name>.txt    what it printed, made by:  sh run_examples.sh
% =====================================================================

\documentclass[aspectratio=169,10pt,compress]{beamer}

% ~~~~~~~~~~~~~~~~~~~~~~~~~~~~~~~~~~~~~~~~~~~~~~~~~~~~~~~~~~~~~~~~~ V hblib
	\usepackage{../../hblib/hblibPresentation} % lib presentation: theme, i/N, \hSection, algorithm2e, ...
	\usepackage{../../hblib/hblibCodeoutput} % lib code + output: \lstcode, \lstcodedown, ...
	\usepackage{../../hblib/hblibDatastructures} % lib data structures: \hString, \hStack, \hQueue
% ~~~~~~~~~~~~~~~~~~~~~~~~~~~~~~~~~~~~~~~~~~~~~~~~~~~~~~~~~~~~~~~~~ A hblib


% xcolor, array (and the L{width} column), listings, tikz and algorithm2e come from the hblib libraries
\usepackage[T1]{fontenc}
\usepackage{lmodern}
\usepackage{booktabs}

\definecolor{gitred}{RGB}{240,80,50}
\newcommand{\cmd}[1]{\texttt{\color{gitred}#1}}

% ---------------------------------------------------------------------
\title{Java on the Command Line}
\subtitle{Sophomore Level: Projects, Classpath, JARs and Build Scripts}
\author[Bingol]{Haluk O. Bingol}

\institute[FCIS]{
    Faculty of Computer and Information Sciences\\
    Yeditepe University
}
\date{
	{\tiny \hbTimeStamp}
}

\begin{document}




% ~~~~~~~~~~~~~~~~~~~~~~~~~~~~~~~~~~~~~~ V frm
\begin{frame}[t,fragile]
  \titlepage
\end{frame}
% ~~~~~~~~~~~~~~~~~~~~~~~~~~~~~~~~~~~~~~ A frm




% ~~~~~~~~~~~~~~~~~~~~~~~~~~~~~~~~~~~~~~ V frm
\begin{frame}[t,fragile] \frametitle{Outline}
  \tableofcontents[hideallsubsections]
\end{frame}
% ~~~~~~~~~~~~~~~~~~~~~~~~~~~~~~~~~~~~~~ A frm

% ~~~~~~~~~~~~~~~~~~~~~~~~~~~~~~~~~~~~~~ V
\hSection[Project]{A Project with Packages}{}{}
% ~~~~~~~~~~~~~~~~~~~~~~~~~~~~~~~~~~~~~~ A

% ~~~~~~~~~~~~~~~~~~~~~~~~~~~~~~~~~~~~~~ V
\subsection[Layout]{Project layout}
% ~~~~~~~~~~~~~~~~~~~~~~~~~~~~~~~~~~~~~~ A




% ~~~~~~~~~~~~~~~~~~~~~~~~~~~~~~~~~~~~~~ V frm
\begin{frame}[t,fragile] \frametitle{A typical project layout}
  \begin{columns}[T]
    \begin{column}{0.48\textwidth}


% +++++++++++++++++++++++++++++++++++++++ V lst
%: -Lst.
\begin{lstlisting}[style=codesnippet, language={}, basicstyle=\ttfamily\scriptsize]
bank/
|-- src/                     sources
|   `-- edu/yeditepe/bank/
|       |-- Main.java
|       |-- Account.java
|       `-- util/Money.java
|-- res/                     resources
|-- lib/                     libraries (jars)
|-- out/                     compiled classes
|-- docs/                    javadoc
|-- build.sh, build.bat      build scripts
`-- app.jar                  the product
\end{lstlisting}
% +++++++++++++++++++++++++++++++++++++++ A lst



    \end{column}
    \begin{column}{0.48\textwidth}
      \begin{itemize}\small
        \item Sources and outputs are \textbf{separate}: delete \texttt{out/} any time, rebuild
        \item The folder path below \texttt{src/} is the \textbf{package} name
        \item Only \texttt{src/}, \texttt{res/} and the scripts go into version control (Git); \texttt{out/}, \texttt{docs/}, \texttt{*.jar} are generated
        \item Maven and Gradle use the same idea with \texttt{src/main/java}, \texttt{src/main/resources}, \texttt{target/} or \texttt{build/}
      \end{itemize}
    \end{column}
  \end{columns}
\end{frame}
% ~~~~~~~~~~~~~~~~~~~~~~~~~~~~~~~~~~~~~~ A frm




% ~~~~~~~~~~~~~~~~~~~~~~~~~~~~~~~~~~~~~~ V frm
\begin{frame}[t,fragile] \frametitle{Compiling many files: an argument file}
  \begin{columns}[T]
    \begin{column}{0.46\textwidth}
      \lstcode[basicstyle=\ttfamily\tiny]{codeoutput/proj/Layout/src/edu/yeditepe/bank/Main.java}{Java}
    \end{column}
    \begin{column}{0.52\textwidth}
      \uncover<2->{\lstoutput[basicstyle=\ttfamily\tiny]{codeoutput/Layout.txt}}
    \end{column}
  \end{columns}

  \vspace{0.2em}
  \small \texttt{@sources.txt} makes \texttt{javac} read its arguments from a file. Windows:
  \cmd{dir /s /b src\textbackslash *.java > sources.txt}. Sub-packages become sub-folders of \texttt{out}.
\end{frame}
% ~~~~~~~~~~~~~~~~~~~~~~~~~~~~~~~~~~~~~~ A frm




% ~~~~~~~~~~~~~~~~~~~~~~~~~~~~~~~~~~~~~~ V frm
\begin{frame}[t,fragile] \frametitle{Letting \texttt{javac} find the sources: \texttt{--source-path}}
  \begin{columns}[T]
    \begin{column}{0.46\textwidth}
      \lstcode[basicstyle=\ttfamily\tiny]{codeoutput/proj/SourcePath/src/edu/yeditepe/bank/Account.java}{Java}
    \end{column}
    \begin{column}{0.52\textwidth}
      \uncover<2->{\lstoutput[basicstyle=\ttfamily\tiny]{codeoutput/SourcePath.txt}}
    \end{column}
  \end{columns}

  \vspace{0.2em}
  \small Given only \texttt{Main.java}, \texttt{javac} compiles every class it needs from \texttt{--source-path}.
  Classes nobody uses are \textbf{not} compiled.
\end{frame}
% ~~~~~~~~~~~~~~~~~~~~~~~~~~~~~~~~~~~~~~ A frm

% ~~~~~~~~~~~~~~~~~~~~~~~~~~~~~~~~~~~~~~ V
\hSection[Classpath]{Libraries and the Classpath}{}{}
% ~~~~~~~~~~~~~~~~~~~~~~~~~~~~~~~~~~~~~~ A

% ~~~~~~~~~~~~~~~~~~~~~~~~~~~~~~~~~~~~~~ V
\subsection[LibJar]{A library jar}
% ~~~~~~~~~~~~~~~~~~~~~~~~~~~~~~~~~~~~~~ A




% ~~~~~~~~~~~~~~~~~~~~~~~~~~~~~~~~~~~~~~ V frm
\begin{frame}[t,fragile] \frametitle{Making and using a library jar}
  \begin{columns}[T]
    \begin{column}{0.42\textwidth}
      \lstcode[basicstyle=\ttfamily\tiny]{codeoutput/proj/LibJar/lib-src/textutil/Text.java}{Java}
      \lstcode[basicstyle=\ttfamily\tiny]{codeoutput/proj/LibJar/src/app/Main.java}{Java}
    \end{column}
    \begin{column}{0.56\textwidth}
      \uncover<2->{\lstoutput[basicstyle=\ttfamily\tiny]{codeoutput/LibJar.txt}}
    \end{column}
  \end{columns}

  \vspace{0.2em}
  \small Three ways to give the classpath: \texttt{-cp} with each jar, the wildcard \texttt{"lib/*"} (all jars in \texttt{lib};
  quote it), or the \texttt{CLASSPATH} environment variable (avoid: it affects every Java program).
\end{frame}
% ~~~~~~~~~~~~~~~~~~~~~~~~~~~~~~~~~~~~~~ A frm




% ~~~~~~~~~~~~~~~~~~~~~~~~~~~~~~~~~~~~~~ V frm
\begin{frame}[t,fragile] \frametitle{The classpath on Windows and macOS}
  \small
  \begin{tabular}{@{}L{0.24\textwidth}L{0.33\textwidth}L{0.33\textwidth}@{}}
    \toprule
                         & \textbf{Windows (cmd)}                         & \textbf{macOS (zsh)} \\
    \midrule
    separator            & \texttt{;}                                     & \texttt{:} \\
    two entries          & \texttt{-cp out;lib\textbackslash a.jar}       & \texttt{-cp out:lib/a.jar} \\
    all jars in a folder & \texttt{-cp "out;lib\textbackslash *"}         & \texttt{-cp "out:lib/*"} \\
    environment variable & \texttt{set CLASSPATH=out;lib\textbackslash a.jar} & \texttt{export CLASSPATH=out:lib/a.jar} \\
    show it              & \texttt{echo \%CLASSPATH\%}                    & \texttt{echo \$CLASSPATH} \\
    \bottomrule
  \end{tabular}

  \vspace{0.5em}
  \begin{itemize}\small
    \item \texttt{-cp} \textbf{replaces} \texttt{CLASSPATH}; without both, the classpath is \texttt{.}
    \item \texttt{lib/*} means all \texttt{.jar} files in \texttt{lib}, not sub-folders and not \texttt{.class} files
    \item In PowerShell, quote arguments with \texttt{;} or \texttt{*}: \texttt{java -cp "out;lib\textbackslash *" app.Main}
  \end{itemize}
\end{frame}
% ~~~~~~~~~~~~~~~~~~~~~~~~~~~~~~~~~~~~~~ A frm

% ~~~~~~~~~~~~~~~~~~~~~~~~~~~~~~~~~~~~~~ V
\hSection[Jar]{JAR Files in Depth}{}{}
% ~~~~~~~~~~~~~~~~~~~~~~~~~~~~~~~~~~~~~~ A

% ~~~~~~~~~~~~~~~~~~~~~~~~~~~~~~~~~~~~~~ V
\subsection[Exec]{Executable jars with libraries}
% ~~~~~~~~~~~~~~~~~~~~~~~~~~~~~~~~~~~~~~ A




% ~~~~~~~~~~~~~~~~~~~~~~~~~~~~~~~~~~~~~~ V frm
\begin{frame}[t,fragile] \frametitle{An executable jar that uses a library}
  \begin{columns}[T]
    \begin{column}{0.3\textwidth}
      \lstcode[basicstyle=\ttfamily\tiny]{codeoutput/proj/ExecJar/manifest.txt}{{}}
    \end{column}
    \begin{column}{0.68\textwidth}
      \uncover<2->{\lstoutput[basicstyle=\ttfamily\tiny]{codeoutput/ExecJar.txt}}
    \end{column}
  \end{columns}

  \vspace{0.2em}
  \small \texttt{Class-Path} in the manifest lists jars \textbf{relative to the jar's own folder}; with \texttt{-jar},
  \texttt{-cp} and \texttt{CLASSPATH} are ignored. Ship the \texttt{dist/} folder as it is.
\end{frame}
% ~~~~~~~~~~~~~~~~~~~~~~~~~~~~~~~~~~~~~~ A frm

% ~~~~~~~~~~~~~~~~~~~~~~~~~~~~~~~~~~~~~~ V
\subsection[Resources]{Resources}
% ~~~~~~~~~~~~~~~~~~~~~~~~~~~~~~~~~~~~~~ A




% ~~~~~~~~~~~~~~~~~~~~~~~~~~~~~~~~~~~~~~ V frm
\begin{frame}[t,fragile] \frametitle{Resources inside a jar}
  \begin{columns}[T]
    \begin{column}{0.66\textwidth}
      \lstcode[basicstyle=\ttfamily\tiny]{codeoutput/proj/Resources/src/app/Hello.java}{Java}
    \end{column}
    \begin{column}{0.32\textwidth}
      \uncover<2->{\lstoutput[basicstyle=\ttfamily\tiny]{codeoutput/Resources.txt}}
    \end{column}
  \end{columns}

  \vspace{0.2em}
  \small Resources (texts, images, settings) are found on the classpath like classes, so they work the same
  from \texttt{out/} and from a jar. Never use file paths such as \texttt{res/app/greeting.txt} in the code.
\end{frame}
% ~~~~~~~~~~~~~~~~~~~~~~~~~~~~~~~~~~~~~~ A frm

% ~~~~~~~~~~~~~~~~~~~~~~~~~~~~~~~~~~~~~~ V
\subsection[Tool]{The jar tool}
% ~~~~~~~~~~~~~~~~~~~~~~~~~~~~~~~~~~~~~~ A




% ~~~~~~~~~~~~~~~~~~~~~~~~~~~~~~~~~~~~~~ V frm
\begin{frame}[t,fragile] \frametitle{The \texttt{jar} tool: create, update, list, extract}
  \lstoutput[basicstyle=\ttfamily\tiny]{codeoutput/JarTool.txt}

  \vspace{0.2em}
  \small Short options like \texttt{tar}: \texttt{c}reate, \texttt{t} list, \texttt{x} extract, \texttt{u}pdate,
  \texttt{f} file, \texttt{e} entry point (main class). \texttt{jar cfe} = \texttt{--create --file --main-class}.
\end{frame}
% ~~~~~~~~~~~~~~~~~~~~~~~~~~~~~~~~~~~~~~ A frm

% ~~~~~~~~~~~~~~~~~~~~~~~~~~~~~~~~~~~~~~ V
\hSection[Docs]{Documentation}{}{}
% ~~~~~~~~~~~~~~~~~~~~~~~~~~~~~~~~~~~~~~ A

% ~~~~~~~~~~~~~~~~~~~~~~~~~~~~~~~~~~~~~~ V
\subsection[Javadoc]{javadoc}
% ~~~~~~~~~~~~~~~~~~~~~~~~~~~~~~~~~~~~~~ A




% ~~~~~~~~~~~~~~~~~~~~~~~~~~~~~~~~~~~~~~ V frm
\begin{frame}[t,fragile] \frametitle{API documentation with \texttt{javadoc}}
  \begin{columns}[T]
    \begin{column}{0.5\textwidth}
      \lstcode[basicstyle=\ttfamily\tiny]{codeoutput/proj/Javadoc/src/edu/yeditepe/bank/util/Money.java}{Java}
    \end{column}
    \begin{column}{0.48\textwidth}
      \uncover<2->{\lstoutput[basicstyle=\ttfamily\tiny]{codeoutput/Javadoc.txt}}
    \end{column}
  \end{columns}

  \vspace{0.2em}
  \small \texttt{/** ... */} comments become HTML pages; open \texttt{docs/index.html} in a browser.
  \texttt{-quiet} hides progress messages. Without \texttt{-Xdoclint:none}, \texttt{javadoc} also warns about every missing
  \texttt{@param}, \texttt{@return} and comment (7 warnings here): useful before a release.
\end{frame}
% ~~~~~~~~~~~~~~~~~~~~~~~~~~~~~~~~~~~~~~ A frm

% ~~~~~~~~~~~~~~~~~~~~~~~~~~~~~~~~~~~~~~ V
\hSection[Build]{Build Scripts}{}{}
% ~~~~~~~~~~~~~~~~~~~~~~~~~~~~~~~~~~~~~~ A

% ~~~~~~~~~~~~~~~~~~~~~~~~~~~~~~~~~~~~~~ V
\subsection[Scripts]{One command to build}
% ~~~~~~~~~~~~~~~~~~~~~~~~~~~~~~~~~~~~~~ A




% ~~~~~~~~~~~~~~~~~~~~~~~~~~~~~~~~~~~~~~ V frm
\begin{frame}[t,fragile] \frametitle{One command to build: \texttt{build.sh} and \texttt{build.bat}}
  \begin{columns}[T]
    \begin{column}{0.49\textwidth}
      \lstcode[basicstyle=\ttfamily\tiny]{codeoutput/proj/Build/build.sh}{bash}
      \small macOS: \cmd{sh build.sh} (or \cmd{chmod +x build.sh} once, then \cmd{./build.sh})
    \end{column}
    \begin{column}{0.49\textwidth}
      \lstcode[basicstyle=\ttfamily\tiny]{codeoutput/proj/Build/build.bat}{{}}
      \small Windows: \cmd{build.bat} in Command Prompt. \texttt{|| exit /b 1} stops on a compile error
      (\texttt{set -e} does that in \texttt{sh}).
    \end{column}
  \end{columns}
\end{frame}
% ~~~~~~~~~~~~~~~~~~~~~~~~~~~~~~~~~~~~~~ A frm




% ~~~~~~~~~~~~~~~~~~~~~~~~~~~~~~~~~~~~~~ V frm
\begin{frame}[t,fragile] \frametitle{Running the build script}
  \lstoutput[basicstyle=\ttfamily\tiny]{codeoutput/Build.txt}

  \vspace{0.2em}
  \small The second run uses a broken \texttt{Account.java}: \texttt{javac} fails, \texttt{set -e} stops the script, and
  its exit code (1) tells the caller (a person, an IDE, or a CI server) that the build failed.
\end{frame}
% ~~~~~~~~~~~~~~~~~~~~~~~~~~~~~~~~~~~~~~ A frm

% ~~~~~~~~~~~~~~~~~~~~~~~~~~~~~~~~~~~~~~ V
\hSection[Config]{Configuring a Program}{}{}
% ~~~~~~~~~~~~~~~~~~~~~~~~~~~~~~~~~~~~~~ A

% ~~~~~~~~~~~~~~~~~~~~~~~~~~~~~~~~~~~~~~ V
\subsection[Config]{Arguments, properties, environment}
% ~~~~~~~~~~~~~~~~~~~~~~~~~~~~~~~~~~~~~~ A




% ~~~~~~~~~~~~~~~~~~~~~~~~~~~~~~~~~~~~~~ V frm
\begin{frame}[t,fragile] \frametitle{Arguments, properties, environment, defaults}
  \begin{columns}[T]
    \begin{column}{0.72\textwidth}
      \lstcode[basicstyle=\ttfamily\tiny]{codeoutput/proj/Config/Config.java}{Java}
    \end{column}
    \begin{column}{0.26\textwidth}
      \uncover<2->{\lstoutput[basicstyle=\ttfamily\tiny]{codeoutput/Config.txt}}
    \end{column}
  \end{columns}

  \vspace{0.2em}
  \small A common order: \textbf{argument} $>$ \textbf{system property} (\texttt{-D}) $>$ \textbf{environment variable} $>$
  \textbf{default}. Windows: \texttt{set PORT=9000} then \texttt{java Config}.
\end{frame}
% ~~~~~~~~~~~~~~~~~~~~~~~~~~~~~~~~~~~~~~ A frm

% ~~~~~~~~~~~~~~~~~~~~~~~~~~~~~~~~~~~~~~ V
\hSection[Summary]{Summary}{}{}
% ~~~~~~~~~~~~~~~~~~~~~~~~~~~~~~~~~~~~~~ A

% ~~~~~~~~~~~~~~~~~~~~~~~~~~~~~~~~~~~~~~ V
\subsection[Cheat sheet]{Cheat sheet}
% ~~~~~~~~~~~~~~~~~~~~~~~~~~~~~~~~~~~~~~ A




% ~~~~~~~~~~~~~~~~~~~~~~~~~~~~~~~~~~~~~~ V frm
\begin{frame}[t,fragile] \frametitle{Cheat sheet}
  \footnotesize
  \begin{tabular}{@{}L{0.5\textwidth}L{0.42\textwidth}@{}}
    \toprule
    \textbf{Command} & \textbf{Does} \\
    \midrule
    \texttt{javac -d out @sources.txt}                          & compile a list of files \\
    \texttt{javac -d out --source-path src src/a/Main.java}     & compile \texttt{Main} and what it needs \\
    \texttt{javac -cp lib/x.jar ...}                            & compile against a library \\
    \texttt{java -cp "out:lib/*" a.Main}                        & run with all jars in \texttt{lib} (Windows \texttt{;}) \\
    \texttt{jar cfe app.jar a.Main -C out .}                    & executable jar \\
    \texttt{jar --create ... --manifest m.txt}                  & add \texttt{Class-Path} etc.\ to the manifest \\
    \texttt{jar tf}, \texttt{jar xf}, \texttt{jar uf}           & list, extract, update \\
    \texttt{javadoc -d docs --source-path src pkg...}           & HTML documentation \\
    \texttt{java -Dkey=value ...}                               & system property \\
    \bottomrule
  \end{tabular}
\end{frame}
% ~~~~~~~~~~~~~~~~~~~~~~~~~~~~~~~~~~~~~~ A frm

% ~~~~~~~~~~~~~~~~~~~~~~~~~~~~~~~~~~~~~~ V
\subsection[Exercises]{Exercises}
% ~~~~~~~~~~~~~~~~~~~~~~~~~~~~~~~~~~~~~~ A




% ~~~~~~~~~~~~~~~~~~~~~~~~~~~~~~~~~~~~~~ V frm
\begin{frame}[t,fragile] \frametitle{Exercises}
  \begin{enumerate}\small
    \item Add a class \texttt{edu.yeditepe.bank.Bank} holding many accounts; rebuild with \texttt{build.sh} \emph{and} \texttt{build.bat}.
    \item Put \texttt{Money} into its own jar \texttt{bank-util.jar}; build \texttt{app.jar} with a \texttt{Class-Path} to it.
    \item Load the currency symbol (\texttt{TL}) from a resource \texttt{money.properties} with \texttt{java.util.Properties}.
    \item Extend \texttt{build.sh} with a \texttt{clean} and a \texttt{docs} target (\texttt{sh build.sh docs}).
    \item What happens with \texttt{java -cp lib/textutil.jar -jar app.jar}? Explain.
    \item Unzip \texttt{app.jar} with your operating system's zip tool. What is inside?
  \end{enumerate}
\end{frame}
% ~~~~~~~~~~~~~~~~~~~~~~~~~~~~~~~~~~~~~~ A frm




% ~~~~~~~~~~~~~~~~~~~~~~~~~~~~~~~~~~~~~~ V frm
\begin{frame}[t,fragile]
	\begin{center}
		\vspace{4cm}
		\hRemark{\Huge Questions?}
	\end{center}
\end{frame}
% ~~~~~~~~~~~~~~~~~~~~~~~~~~~~~~~~~~~~~~ A frm



\end{document}
